\section{\texorpdfstring{近似 $|\bnabla\psi| \approx \delta_s$ の精度：Eikonal 条件下で $\Ord{\varepsilon^2}$}{近似 |∇ψ|≈δ\_s の精度：Eikonal 条件下で O(ε²)}}
\label{app:csf_delta_precision}
% ═══════════════════════════════════════════════════════════════════════════════

Eikonal 条件 $|\bnabla\phi|=1$ が厳密に成立するとき，
§\ref{sec:csf} の近似 $|\bnabla\psi| \approx \delta_s$ の誤差を
法線方向座標によるテイラー展開で定量化する．

\medskip
\noindent\textbf{法線座標展開}

界面 $\Gamma$（$\phi=0$）の法線方向座標を $s$，曲率 $\kappa_1, \kappa_2$ を持つ界面近傍の体積要素は：
\[
  dV = \bigl(1 - \kappa_1 s\bigr)\bigl(1 - \kappa_2 s\bigr)\, ds\, dA_\Gamma
  \approx \bigl(1 - 2H s + \kappa_1\kappa_2 s^2 + \cdots\bigr)\, ds\, dA_\Gamma
\]
ここで $H = (\kappa_1+\kappa_2)/2$ は平均曲率．
界面デルタ関数 $\delta_s$ は $\int_\Omega f\,\delta_s\,dV = \int_\Gamma f\,dA_\Gamma$ を満たす分布として
定義される．正則化 $\delta_\varepsilon(s)$ を用いた体積積分では，Jacobian 因子
$(1 - 2Hs + \kappa_1\kappa_2 s^2 + \cdots)$ が被積分関数に現れる（下記参照）．

\noindent\textbf{近似誤差の評価}

$|\bnabla\phi|=1$ のとき $|\bnabla\psi| = \delta_\varepsilon(s)$（式 \eqref{eq:delta_eps_def}）であるから：
\begin{align*}
  \int_\Omega f\,|\bnabla\psi|\,dV - \int_\Gamma f\,dA_\Gamma
  &= \int_\Gamma \int_{-\infty}^{\infty}
       f\,\delta_\varepsilon(s)\,\bigl(1 - 2Hs + \Ord{s^2}\bigr)\,ds\, dA_\Gamma
     - \int_\Gamma f\, dA_\Gamma \\
  &= \int_\Gamma f\!\int_{-\infty}^{\infty}\delta_\varepsilon(s)\cdot(-2Hs)\,ds\, dA_\Gamma
     + \Ord{\varepsilon^2}
\end{align*}
$\delta_\varepsilon(s)$（式 \eqref{eq:delta_eps_def}）は $s$ の\textbf{偶関数}であり，$s$ は奇関数であるから：
\[
  \int_{-\infty}^{\infty} s\,\delta_\varepsilon(s)\,ds = 0
\]
したがって曲率 $H$ の項は消え，誤差の主要項は $s^2\,\delta_\varepsilon(s)$ の積分から来る．
置換 $t = s/\varepsilon$ により：
\[
  \int_{-\infty}^{\infty} s^2\,\delta_\varepsilon(s)\,ds
  = \varepsilon^2 \int_{-\infty}^{\infty} t^2\,\frac{e^{-t}}{(1+e^{-t})^2}\,dt
  = \frac{\pi^2\varepsilon^2}{3}
\]
最後の等号は次のように得られる．被積分関数を級数展開すると
$t^2 e^{-t}/(1+e^{-t})^2 = t^2 \sum_{n=1}^{\infty}(-1)^{n+1} n\,e^{-nt}$．
$[0,\infty)$ 上で項別積分し $\int_0^\infty t^2 n\,e^{-nt}\,dt = 2/n^2$ を用いると
$\int_0^\infty (\cdot)\,dt = 2\sum_{n=1}^{\infty}(-1)^{n+1}/n^2 = 2\eta(2) = \pi^2/6$
（ディリクレのエータ関数 $\eta(2) = \pi^2/12$）．
$[-\infty,0]$ 側は偶関数性で等値であるから，全体で $\pi^2/3$ を得る．

\noindent\textbf{結論}
\[
  \int_\Omega f\,|\bnabla\psi|\,dV - \int_\Gamma f\,dA_\Gamma
  = \Ord{\varepsilon^2\,\kappa^2\,\|f\|_{L^\infty}}
\]
Eikonal 条件下での近似 $|\bnabla\psi| \approx \delta_s$ の誤差は $\Ord{\varepsilon^2}$ である
（より正確には $\Ord{\varepsilon^2\kappa^2}$；界面が直線 $\kappa=0$ のときは厳密）．
$\varepsilon \sim \Delta x$ の設定では誤差は $\Ord{\Delta x^2}$．\qed

% ═══════════════════════════════════════════════════════════════════════════════
